% ---------------------------------------------------------------------
%  Chapter 9 : Diagram Enumeration -- Prediagrams, nauty, and the
%              loop-order completeness bound
% ---------------------------------------------------------------------
\chapter{Diagram Enumeration}\label{ch:enumeration}

Every loop calculation in this manual ultimately reduces to a sum over Feynman
diagrams. Before the engine can evaluate that sum, it must first \emph{know which
diagrams are in it} --- and it must know that the list is \emph{complete}, with
nothing missing and nothing double-counted. That bookkeeping is the job of the
subsystem dissected in this chapter. Given just two integers --- $k$, the number
of external legs of the correlation function we want, and $\ell$, the loop order
of the correction we are computing --- it produces the complete, deduplicated
list of every bare graph topology that could possibly contribute, \emph{before
any physical field labels are attached}. Those bare graphs are called
\term{prediagrams}.

The whole subsystem lives in two files and we read both in full:
\file{msrjd/enumeration/loop\_diagram\_enumeration.py} (the topological
generator) and \file{msrjd/enumeration/degree\_scan.py} (a small feedback bridge
back to the algebra). The generator is \emph{theory-agnostic}: it never looks at
the action, the fields, or the parameters. It answers a purely combinatorial
question --- \emph{what graph shapes are even possible at $(k,\ell)$?} --- and
leaves the physical question --- \emph{which of those shapes the theory actually
populates, and with what propagators and vertices} --- to the type-assignment
subsystem in the next chapter.

The reader is assumed to know \msrjd{} field theory but \emph{not} the tooling.
The heavy lifting here is done by \term{SageMath} and, inside it, by the graph
canonical-labeling engine \term{nauty}; both are introduced from first
principles, as are the handful of Python standard-library helpers
(\code{itertools}, \code{collections}, \code{concurrent.futures}) the module
leans on. When in doubt we over-explain.

\section{What a prediagram is, and why we strip the physics first}

In \msrjd{} (Martin--Siggia--Rose--Janssen--De~Dominicis) field theory, a genuine
Feynman diagram is a graph in which \emph{every edge is a specific propagator}
and \emph{every internal vertex is a specific interaction monomial} from the
action. The doubled field content of \msrjd{} means there are two flavors of
propagator: a \emph{response} propagator $\avg{\phi\,\tilde\phi}$ (causal,
directed --- it carries a definite time direction, from cause to effect) and a
\emph{correlation} propagator $\avg{\phi\,\phi}$. The vertices come from the
Taylor-expanded interaction part of the action --- a cubic vertex $\tilde\phi\,
\phi^2$, a noise vertex $\tilde\phi^2$, and so on.

A \term{prediagram} strips all of that physics away and keeps only the
combinatorial skeleton. Concretely it retains exactly three things:

\begin{itemize}
  \item the \textbf{shape} of the graph --- which vertices connect to which;
  \item the \textbf{leaf-vs-internal} distinction --- which vertices are
        \emph{external legs} of the correlation function, and which are
        \emph{internal} interaction points; and
  \item an \textbf{orientation} on every edge (an arrow $u \to v$), because
        \msrjd{} propagators are causal: a response propagator flows in a
        definite time direction, so the edges of a valid diagram must carry
        arrows.
\end{itemize}

So a prediagram is, in one phrase, ``a Feynman diagram with the physics not yet
filled in.''

\begin{defn}
A \textbf{prediagram} is a bare, oriented graph topology: a partition of vertices
into \emph{leaves} (external legs) and \emph{internal} vertices, together with an
arrow on every edge, but with \emph{no} propagator types and \emph{no} vertex
monomials assigned yet. It is the output object of the enumeration subsystem,
carried as the 4-tuple \code{(D, G, leaves, internal)}.
\end{defn}

\subsection{Why separate the topology from the physics?}

Two reasons, one practical and one conceptual.

The conceptual reason is that the set of admissible \emph{shapes} depends only on
$k$ and $\ell$ --- it is the same for an Ornstein--Uhlenbeck particle, the
Kardar--Parisi--Zhang equation, or a coupled multi-field reaction--diffusion
model. Solving the topology problem once, model-independently, means the
expensive graph-isomorphism work is never repeated per theory.

The practical reason is that it cleanly factors a hard combinatorial problem into
two tractable stages. Stage one (this chapter) enumerates shapes. Stage two
(Chapter~\ref{ch:type-assignment}) walks each shape and tries every consistent
way to label its edges with propagators and its vertices with monomials, throwing
away the labelings the theory does not support. Keeping the two apart means each
stage has a single, checkable correctness criterion: \emph{completeness} for the
topology stage (we miss no shape), and \emph{validity} for the typing stage (we
keep only labelings the action provides).

\subsection{Where it sits in the pipeline}

The project is organized into lettered ``build phases'' (documented in
\file{docs/BUILD\_PHASE\_OUTLINES.md}). This subsystem is the topological
generator --- it is invoked as phase~\code{[5/7]} of the temporal trace you saw
in Chapter~\ref{ch:quickstart} (``Enumerate prediagrams''). Schematically:

\begin{lstlisting}[style=console]
  Phase 1 (FieldTheory.expand)     Phase 2 (THIS SUBSYSTEM)        Phase D/E
  ----------------------------     -------------------------       ------------
  Action -> Taylor-expanded        enumerate_prediagrams(k,ell)    type_assignment
  vertices + noise kernel,     --> trees -> topologies ->      --> .enumerate_typed
  saved as a "theory"              prediagrams                     _diagrams
                                        |
                                        v
                               Phase C: degree_scan
                               .ensure_taylor_order
                               (feedback to Phase 1)
\end{lstlisting}

\begin{itemize}
  \item \textbf{What feeds it:} essentially nothing but the two integers $k$ and
        $\ell$. The generation is theory-agnostic; the action enters only later.
  \item \textbf{What consumes it:} two downstream readers.
        \code{msrjd/diagrams/type\_assignment.py} is the direct consumer --- its
        \code{enumerate\_typed\_diagrams(prediagram, ...)} unpacks one prediagram
        tuple \code{(D, G, leaves, internal)} and produces every valid fully-typed
        diagram for it. And \code{msrjd/enumeration/degree\_scan.py} (build
        phase~C) scans the same tuples for their maximum vertex degree, to decide
        whether the saved theory was Taylor-expanded far enough --- the feedback
        bridge of \S\ref{sec:degree-scan}.
\end{itemize}

\subsection{The three-stage refinement}

The generator works in three successive refinements. Each stage is a
\emph{filter-and-deduplicate} pass on the previous one:

\begin{enumerate}
  \item \textbf{Trees} --- connected acyclic graphs with the right number of
        leaves. A tree is the loop-free ``spanning skeleton'' of every possible
        diagram. Trees have zero loops, so to reach $\ell$ loops we will later add
        $\ell$ extra edges.
  \item \textbf{Topologies} --- undirected multigraphs obtained by adding exactly
        $\ell$ extra edges to a tree (creating exactly $\ell$ independent
        cycles), keeping only those that satisfy \msrjd's structural constraints,
        then removing isomorphic duplicates.
  \item \textbf{Prediagrams} --- topologies with every edge given an arrow,
        keeping only physically admissible (causal, acyclic) orientations, then
        removing isomorphic duplicates again.
\end{enumerate}

The public entry points (\code{enumerate\_all}, \code{enumerate\_prediagrams},
and friends) wrap these three stages; \S\ref{sec:api} is the reference.

\section{A graph-theory primer for field theorists}

This subsystem is pure graph theory, so we build the vocabulary from the ground
up. A reader who knows \msrjd{} but not the combinatorics should be able to
follow every later step.

\subsection{Vertices, edges, degree, leaves}

A \term{graph} is a set of \emph{vertices} (points) and \emph{edges}
(connections between pairs of points). The \term{degree} of a vertex $v$, written
$\deg(v)$, is the number of edge-ends touching it. If two vertices are joined by
a \emph{double edge} --- a multi-edge --- that contributes $2$ to each of their
degrees.

We partition the vertices by degree:

\begin{itemize}
  \item \term{Leaves} are the degree-1 vertices. In our setting these are the
        \emph{external legs} of the correlation function; a $k$-point function has
        exactly $k$ of them in the final diagram.
  \item \term{Internal vertices} are the degree-$\ge 2$ vertices. These are the
        interaction points.
  \item Internal vertices split further into \emph{degree-2} vertices ($V_2$) and
        \emph{degree-$\ge 3$} vertices ($V_3$, written \code{degree\_3plus} in the
        code).
\end{itemize}

In code this partition is computed by \code{classify\_vertices\_sage}
(\file{loop\_diagram\_enumeration.py:32}), which returns the four lists
\code{(leaves, internal, degree\_2, degree\_3plus)} by one filtered
list-comprehension per category over \code{G.vertices()}:

\begin{lstlisting}
def classify_vertices_sage(G):
    leaves       = [v for v in G.vertices() if G.degree(v) == 1]
    internal     = [v for v in G.vertices() if G.degree(v) > 1]
    degree_2     = [v for v in G.vertices() if G.degree(v) == 2]
    degree_3plus = [v for v in G.vertices() if G.degree(v) >= 3]
    return leaves, internal, degree_2, degree_3plus
\end{lstlisting}

Note \code{internal} is the union of \code{degree\_2} and \code{degree\_3plus}.
This is the single most-called helper in the module --- nearly every constraint
check begins by classifying the vertices.

\subsection{Trees and the loop number}

A \term{tree} is a connected graph with no cycles. The defining counting
identity of a tree is
\begin{equation}
  |E| \;=\; |V| - 1
  \qquad\text{(a tree on $|V|$ vertices has exactly $|V|-1$ edges).}
  \label{eq:tree-edges}
\end{equation}
For \emph{any} connected graph $G$, the number of \term{independent cycles} ---
equivalently the \term{loop number} or first Betti number --- is
\begin{equation}
  L(G) \;=\; |E| - |V| + 1.
  \label{eq:loop-number}
\end{equation}
This is exactly \code{count\_cycles\_sage} (\file{loop\_diagram\_enumeration.py:68}):

\begin{lstlisting}
def count_cycles_sage(G):
    if not G.is_connected():
        return -1
    return G.size() - G.order() + 1
\end{lstlisting}

In SageMath, \code{G.size()} is the number of edges $|E|$ and \code{G.order()} is
the number of vertices $|V|$. The function returns $-1$ as a sentinel when the
graph is disconnected, since \eqref{eq:loop-number} is only meaningful for a
connected graph. (That $-1$ is load-bearing: it is never equal to any $\ell\ge0$,
so a disconnected candidate is automatically rejected by the
``\code{== ell}'' loop test downstream --- see Gotcha~\ref{gotcha:sentinel}.)

\begin{note}
The consequence we use \emph{everywhere}: starting from a tree ($L=0$) and adding
$\ell$ extra edges \emph{without adding new vertices} raises $|E|$ by $\ell$ and
leaves $|V|$ unchanged, so by \eqref{eq:loop-number} it raises the loop number by
exactly $\ell$. This is the engine of the whole topology stage: \emph{take a
tree, add $\ell$ edges, demand $L(G)=\ell$.}
\end{note}

\section{The leaf surplus \texorpdfstring{$j$}{j} and the loop budget}

The final diagram must have exactly $k$ external legs. But the \emph{intermediate
trees} may temporarily carry \emph{more} than $k$ leaves, because some of those
extra leaves will be ``eaten'' when we add the loop edges: an added edge can
attach to a former leaf, raising its degree above 1 so it stops being a leaf.

The code calls this surplus $j$. A tree is generated with $k+j$ leaves, where
$j\ge 0$ counts the leaves that the $\ell$ added edges will consume.

\subsection{How large can the surplus be?}

Each added edge has two endpoints, so $\ell$ added edges provide $2\ell$
endpoint-attachments. In the most leaf-consuming case, attaching those endpoints
to leaves is what converts leaves into internal vertices. The code caps the
surplus at
\begin{equation}
  j_{\max} \;=\; \ell + \left\lfloor \tfrac{\ell}{2} \right\rfloor
  \;=\; \left\lfloor \tfrac{3\ell}{2} \right\rfloor,
  \label{eq:jmax}
\end{equation}
implemented verbatim as \code{j\_max = ell + (ell // 2)}
(\file{loop\_diagram\_enumeration.py:124}; \code{//} is Python's integer
floor-division). Intuitively this is the maximum number of tree-leaves the
$\ell$ loop-edges can absorb across all valid decompositions; the $3\ell/2$ form
is the integer ceiling on how many degree-1 nodes the loop structure can use up.
For $\ell=1$ this gives $j_{\max}=1$; for $\ell=2$, $j_{\max}=3$; for $\ell=3$,
$j_{\max}=4$.

\section{The degree bounds and the completeness theorem}
\label{sec:bounds}

Not every tree with $k+j$ leaves can be the skeleton of a valid \msrjd{} diagram.
Two structural bounds prune the tree search. Both are computed at the top of the
$j$-loop in \code{generate\_trees\_with\_constraints}
(\file{loop\_diagram\_enumeration.py:118}).

\subsection{Bound on degree-\texorpdfstring{$\ge 3$}{>=3} vertices}

The first bound limits the number of high-degree (branching) vertices:
\begin{equation}
  |V_3| \;\le\; k + j - 2 \;=\; \code{v3\_max}
  \qquad (\file{loop\_diagram\_enumeration.py:128}).
  \label{eq:v3max}
\end{equation}
This is the standard fact about trees: a tree with $m$ leaves has at most $m-2$
vertices of degree $\ge 3$ (the extreme case being a binary tree, where every
internal branching is exactly three-way). Here $m = k+j$, so $\code{v3\_max} =
k+j-2$.

\subsection{Bound on degree-2 vertices: the completeness bound}

The second bound is the subtle one --- the \term{enumeration completeness bound}
referenced throughout the project. The code
(\file{loop\_diagram\_enumeration.py:139}) sets
\begin{lstlisting}
v2_max = 2 * v3_max + 3 * ell - k - 3 * j + 3
\end{lstlisting}
and the in-code comment records that this algebraically equals the \emph{proven}
upper bound on the number of degree-2 vertices of the final topology $G$:
\begin{equation}
  |V_2^{T}| \;\le\; k + 3\ell - j - 1.
  \label{eq:v2max}
\end{equation}
Let us verify the algebra by substituting $\code{v3\_max}=k+j-2$:
\begin{align*}
  2\,\code{v3\_max} + 3\ell - k - 3j + 3
   &= 2(k+j-2) + 3\ell - k - 3j + 3 \\
   &= 2k + 2j - 4 + 3\ell - k - 3j + 3 \\
   &= k - j - 1 + 3\ell \;=\; k + 3\ell - j - 1. \qquad\checkmark
\end{align*}
So the code's \code{v2\_max} is exactly the bound \eqref{eq:v2max}.

\subsection{Why \texorpdfstring{$k+3\ell-j-1$}{k+3l-j-1} is the right bound}

The full proof lives in the paper appendix; the in-code comment
(\file{loop\_diagram\_enumeration.py:129--138}) and the project note
\file{project\_enumeration\_bound\_fix.md} summarize the three ingredients:

\begin{enumerate}
  \item a \textbf{degree-partition identity} --- the handshake/incidence count
        relating leaves, $V_2$, $V_3$, edges, and the loop number;
  \item the \textbf{orientability constraint} $|V_2^{G}| \le k + \ell - 1$ --- a
        degree-2 internal vertex of an \msrjd{} diagram must be orientable as
        one-in/one-out, and there is a hard ceiling on how many such vertices a
        valid causal orientation can admit; and
  \item an \textbf{incidence-counting} bound $\theta \le 2\ell - j$, where
        $\theta$ counts certain doubled-edge / multi-edge incidences.
\end{enumerate}
Combining the three yields \eqref{eq:v2max}.

\begin{defn}
The \textbf{completeness bound} $|V_2^{T}| \le k+3\ell-j-1$ is the proven upper
bound on the number of degree-2 vertices a valid topology can have at fixed
$(k,\ell,j)$. Pruning the tree search by this bound is what \emph{guarantees}
the search misses no contributing topology --- this is the sense in which the
enumeration is provably complete.
\end{defn}

\subsection{The bound-correction story: a false lemma at \texorpdfstring{$\ell\ge 3$}{l>=3}}

This bound has a cautionary history worth telling, because it encodes a principle
the whole subsystem now lives by. An \emph{earlier} version of the code used a
\emph{tighter} bound carrying an extra $-\lfloor j/3\rfloor$ term:
\begin{equation}
  \code{v2\_max\_OLD} \;=\; k + 3\ell - j - 1 - \left\lfloor \tfrac{j}{3}
  \right\rfloor .
  \label{eq:v2old}
\end{equation}
That extra tightening came from a claimed lemma $\theta \le 2\ell - j -
\lfloor j/3\rfloor$, which is \emph{false at $\ell\ge 3$}. The counterexample,
recorded in the comment at \file{loop\_diagram\_enumeration.py:133} (and in
\file{scratch/bound\_check.py}), is \emph{three doubled-edge bubbles on a hub},
$|V|=8$, where every decomposition has $j=3$ and $\theta = 3 > 2\ell - j -
\lfloor j/3\rfloor$.

The insidious part: the false lemma never actually \emph{bound} at the small
orders that had been tested. Exhaustive enumeration at $(k,\ell)\in\{(2,1),(3,1),
(2,2),(3,2),(4,1),(2,3)\}$ showed a slack of $\ge 2$, so the tighter
\eqref{eq:v2old} and the true \eqref{eq:v2max} produced \emph{identical output} at
every tested order. Only the proven bound \eqref{eq:v2max}, however,
\emph{guarantees} completeness at $\ell\ge 3$.

\begin{gotcha}
\label{gotcha:proof-not-slack}
The headline lesson, recorded in the project memory: \textbf{a prune must follow
only a proven bound; empirical slack at small orders is not a proof.} If anyone
re-tightens \code{v2\_max} (or \code{j\_max}), the new expression must come from a
\emph{proven} statement. The June~2026 fix (commit \code{40454e7}) removed the
false $-\lfloor j/3\rfloor$ term; it changed only the \emph{proof}, not the
observed output at the tested orders. The deduplicated \emph{topology} counts at the tested orders ---
$(k,\ell)=(2,1)\to 9$, $(3,1)\to 67$, $(2,2)\to 289$ --- are the same before and
after the fix. (The comment at \file{loop\_diagram\_enumeration.py:138} records
the \emph{orders} that were exhaustively slack-verified,
$\{(2,1),(3,1),(2,2),(3,2),(4,1),(2,3)\}$, not these counts.)
\end{gotcha}

\section{Stage 1: generating the trees}

With the bounds in hand we can read the tree generator. The function
\code{generate\_trees\_with\_constraints(k, ell, max\_vertices\_search=50)}
(\file{loop\_diagram\_enumeration.py:118}) returns a list of triples
\code{(tree, j, num\_leaves)} --- every non-isomorphic tree satisfying the degree
constraints, tagged with its surplus $j$ and its leaf count $\code{num\_leaves} =
k+j$.

\begin{lstlisting}
def generate_trees_with_constraints(k, ell, max_vertices_search=50):
    valid_trees = []
    j_max = ell + (ell // 2)

    for j in range(0, j_max + 1):
        num_leaves = k + j
        v3_max = k + j - 2
        v2_max = 2 * v3_max + 3 * ell - k - 3 * j + 3

        min_n = num_leaves if num_leaves == 1 else num_leaves + 1
        max_n = min(num_leaves + v2_max + v3_max, max_vertices_search)
        max_n = min(max_n, num_leaves + min(10, v2_max + v3_max))

        if max_n < min_n:
            continue

        for n in range(min_n, max_n + 1):
            for tree in graphs.trees(n):
                leaves, internal, degree_2, degree_3plus = classify_vertices_sage(tree)
                if len(leaves) != num_leaves:
                    continue
                if len(degree_3plus) > v3_max:
                    continue
                if len(degree_2) > v2_max:
                    continue
                valid_trees.append((tree, j, num_leaves))

    return valid_trees
\end{lstlisting}

Reading it in order:

\begin{enumerate}
  \item $j_{\max}$ is set by \eqref{eq:jmax}, and we loop $j$ from $0$ to
        $j_{\max}$. For each $j$ we fix $\code{num\_leaves}=k+j$, then compute
        \code{v3\_max} \eqref{eq:v3max} and \code{v2\_max} \eqref{eq:v2max}.
  \item The vertex-count search window $[\code{min\_n}, \code{max\_n}]$ is set:
        \code{min\_n} is \code{num\_leaves} if there is a single leaf (the
        degenerate one-vertex tree), else $\code{num\_leaves}+1$ (you need at
        least one internal vertex). The upper end is $\code{num\_leaves}+
        \code{v2\_max}+\code{v3\_max}$ (the most vertices the bounds permit),
        capped by \code{max\_vertices\_search}, and then capped \emph{again} by
        the \code{+10} guard discussed below.
  \item If the window is empty (\code{max\_n < min\_n}), this $j$ is skipped.
  \item For each $n$ in the window, \code{graphs.trees(n)} yields every
        non-isomorphic tree on $n$ vertices exactly once (\S\ref{sec:sage}). A
        tree is \emph{kept} iff its leaf count equals \code{num\_leaves}, its
        $V_3$ count is $\le\code{v3\_max}$, and its $V_2$ count is
        $\le\code{v2\_max}$.
\end{enumerate}

\begin{gotcha}
\label{gotcha:plus10}
\textbf{The \code{+10} vertex-search cap can in principle undercut completeness at
high order.} The line
\code{max\_n = min(max\_n, num\_leaves + min(10, v2\_max + v3\_max))}
(\file{loop\_diagram\_enumeration.py:143}) clamps the \emph{internal}-vertex
budget to at most $10$. The proven bound only guarantees that all needed
topologies fit within $\code{num\_leaves}+\code{v2\_max}+\code{v3\_max}$
vertices. For small $(k,\ell)$ the proven budget is well under $10$, so the cap is
inert. But at large enough orders $\code{v2\_max}+\code{v3\_max} > 10$, and then
the cap could \emph{skip trees the proven bound says we need}, silently dropping
topologies. This is a performance guard, not a proven-safe bound --- flag it for
review at high order.
\end{gotcha}

\section{Stage 2: adding edges to make topologies}

A tree has $L=0$. To reach $\ell$ loops we add exactly $\ell$ edges, then keep
only the results that have exactly $\ell$ loops, satisfy the structural
constraints, and are not isomorphic duplicates of something already kept.

\subsection{Enumerating the edges to add}

The candidate edges are produced by \code{generate\_edge\_multisets}
(\file{loop\_diagram\_enumeration.py:166}):

\begin{lstlisting}
def generate_edge_multisets(all_vertices, ell):
    if len(all_vertices) < 2:
        return []
    possible_edges = list(combinations(all_vertices, 2))
    return list(combinations_with_replacement(possible_edges, ell))
\end{lstlisting}

The two \code{itertools} generators here matter (\S\ref{sec:itertools}):
\code{combinations(all\_vertices, 2)} lists every \emph{possible} edge --- every
unordered pair of distinct vertices. Then
\code{combinations\_with\_replacement(possible\_edges, ell)} lists every
\emph{multiset} of $\ell$ edges to add. ``With replacement'' is essential: it
permits choosing the \emph{same} edge twice, which creates a double-edge between
that pair of vertices --- a \emph{bubble}, the simplest one-loop topology. A
plain \code{combinations} would forbid bubbles and the enumeration would be
incomplete.

The chosen edges are spliced onto a copy of the tree by \code{add\_edges\_to\_tree}
(\file{loop\_diagram\_enumeration.py:173}), which builds the copy as a Sage
\code{Graph(tree, multiedges=True, loops=False)} --- multi-edges \emph{allowed}
(we just argued we need them), self-loops \emph{forbidden} (an edge from a vertex
to itself is never physical here).

\subsection{The structural constraints on a topology}

Beyond having the right loop number, a topology must satisfy local structural
rules so that it reduces to a unique canonical diagram. These are enforced by
\code{check\_topology\_constraints(G, k, ell=None)}
(\file{loop\_diagram\_enumeration.py:180}):

\begin{lstlisting}
def check_topology_constraints(G, k, ell=None):
    leaves, internal, degree_2, degree_3plus = classify_vertices_sage(G)
    if len(leaves) != k:
        return False
    if not G.is_connected():
        return False
    if ell is None or ell > 0:
        if has_adjacent_degree2_sage(G, degree_2):
            return False
        if not check_deg3_has_non_deg2_neighbor(G, degree_3plus, degree_2):
            return False
        if not check_leaf_neighbors_not_all_deg2(G, leaves, degree_2):
            return False
    return True
\end{lstlisting}

The unconditional checks are the obvious ones: the final topology must have
exactly $k$ external legs (\code{len(leaves) == k}), and it must be
\code{is\_connected} (a diagram contributing to a \emph{connected} cumulant must
be connected). The three conditional checks are all ``degree-2 pruning'' rules
that remove redundant subdivisions of a single propagator line:

\begin{description}
  \item[No two adjacent degree-2 vertices] (\code{has\_adjacent\_degree2\_sage},
        \file{loop\_diagram\_enumeration.py:41}). A chain of two degree-2 vertices
        in a row is a redundant subdivision of one propagator line; only one
        canonical representative is kept.
  \item[Every degree-$\ge3$ vertex has a non-degree-2 neighbor]
        (\code{check\_deg3\_has\_non\_deg2\_neighbor},
        \file{loop\_diagram\_enumeration.py:50}). This returns \code{False} if
        some branching vertex is surrounded \emph{only} by subdivision chains ---
        an ``isolated hub.''
  \item[Leaf neighbors are not all degree-2]
        (\code{check\_leaf\_neighbors\_not\_all\_deg2},
        \file{loop\_diagram\_enumeration.py:58}). An external leg should not
        attach \emph{only} through a chain of degree-2 subdivisions.
\end{description}

\begin{gotcha}
\label{gotcha:ell0-skip}
The three degree-2 pruning checks are \emph{skipped at tree level} ($\ell=0$). The
gate is \code{if ell is None or ell > 0:} (\file{loop\_diagram\_enumeration.py:189}),
and the comment explains why: at $\ell=0$ ``the tree IS the topology --- no
contraction ambiguity'' (\file{loop\_diagram\_enumeration.py:188}). The constraint
set is therefore \emph{different} at tree level vs.\ loop level. Passing the wrong
\code{ell} into this function would change which topologies survive --- a subtle
correctness trap.
\end{gotcha}

\subsection{Processing one tree, then deduplicating}

The unit of work is \code{process\_tree\_parallel(args)}
(\file{loop\_diagram\_enumeration.py:203}), where \code{args = (tree, j,
num\_leaves, k, ell)}. For each edge-multiset it adds the edges, demands
\code{count\_cycles\_sage(G) == ell}, demands \code{check\_topology\_constraints},
then relabels the survivor leaves-first and re-classifies:

\begin{lstlisting}
def process_tree_parallel(args):
    tree, j, num_leaves, k, ell = args
    local_candidates = []
    all_verts = list(tree.vertices())
    for edge_multiset in generate_edge_multisets(all_verts, ell):
        G = add_edges_to_tree(tree, edge_multiset)
        if count_cycles_sage(G) != ell:
            continue
        if not check_topology_constraints(G, k, ell=ell):
            continue
        G_relabeled, _ = relabel_leaves_first(G)
        leaves_final, internal_final, _, _ = classify_vertices_sage(G_relabeled)
        local_candidates.append((G_relabeled, leaves_final, internal_final))
    return local_candidates
\end{lstlisting}

The name says ``parallel'' because this is the function submitted to the thread
pool (\S\ref{sec:parallel}); it runs perfectly well serially too. The relabeling
step \code{relabel\_leaves\_first} (\file{loop\_diagram\_enumeration.py:74})
renumbers the graph so that \emph{leaves are $0\ldots|L|-1$ and internal vertices
are $|L|\ldots|V|-1$}; it calls \code{G.relabel(dict, inplace=False)}, which
(because of \code{inplace=False}) returns a \emph{new} relabeled graph rather than
mutating the original. This canonical leaf-first labeling is the invariant the
downstream code relies on (Gotcha~\ref{gotcha:leaffirst}).

The candidates from all trees are then collapsed by
\code{\_remove\_isomorphic\_undirected} (\file{loop\_diagram\_enumeration.py:219}):

\begin{lstlisting}
def _remove_isomorphic_undirected(candidates):
    unique = []
    for G, leaves, internal in candidates:
        if not any(graphs_isomorphic_with_labels(G, leaves, G2, l2)
                   for G2, l2, _ in unique):
            unique.append((G, leaves, internal))
    return unique
\end{lstlisting}

This is a classic $O(n^2)$ dedup: walk the candidates, keep one only if it is
\emph{not} isomorphic to any already-kept representative. The isomorphism
predicate \code{graphs\_isomorphic\_with\_labels} is the subject of
\S\ref{sec:iso}; each comparison is cheap because it fast-rejects on cheap
invariants before invoking the expensive test.

\section{Stage 3: orienting the edges into prediagrams}

An \msrjd{} edge is a \emph{directed} propagator, so each undirected topology must
be oriented. With $|E|$ edges there are $2^{|E|}$ raw orientations; only the
physically admissible ones survive.

\subsection{Building one orientation}

\code{orient\_edges(G, orientation\_bits)}
(\file{loop\_diagram\_enumeration.py:264}) turns a bit-vector into a
\code{DiGraph}:

\begin{lstlisting}
def orient_edges(G, orientation_bits):
    D = DiGraph(multiedges=True, loops=False)
    D.add_vertices(G.vertices())
    for i, (u, v) in enumerate(G.edges(labels=False)):
        if orientation_bits[i] == 0:
            D.add_edge(u, v, i)
        else:
            D.add_edge(v, u, i)
    return D
\end{lstlisting}

The $i$-th edge keeps its direction $u\to v$ when $\code{orientation\_bits}[i]=0$,
and flips to $v\to u$ otherwise. Crucially \textbf{the edge index $i$ is used as
the edge label} (\code{D.add\_edge(u, v, i)}). This preserves multi-edge identity:
when a topology has a double edge between the same pair, the two parallel arrows
remain distinguishable by their labels, and downstream code
(\code{type\_assignment} iterates \code{D.edges()} as 3-tuples \code{(u, v,
label)}) depends on those labels being present and distinct
(Gotcha~\ref{gotcha:edgelabels}). The per-edge order comes from
\code{G.edges(labels=False)}, which is deterministic for a fixed graph object.

\subsection{The MSR-JD causality constraints}

The admissible orientations are selected by \code{check\_orientation\_constraints}
(\file{loop\_diagram\_enumeration.py:275}):

\begin{lstlisting}
def check_orientation_constraints(D, leaves):
    leaves_set = set(leaves)
    if not D.is_directed_acyclic():
        return False
    for v in D.vertices():
        in_deg  = D.in_degree(v)
        out_deg = D.out_degree(v)
        if (in_deg + out_deg) == 1 and in_deg != 1:
            return False
        if out_deg == 0 and in_deg >= 2:
            return False
        if in_deg == 1 and out_deg == 1:
            return False
    sources = [v for v in D.vertices() if D.in_degree(v) == 0]
    for v in sources:
        if any(u in sources for u in D.neighbors_out(v)):
            return False
    return True
\end{lstlisting}

Each clause has a direct physical reading:

\begin{description}
  \item[Acyclic.] \code{D.is\_directed\_acyclic()} --- \msrjd{} response
        propagators are causal, so a valid diagram has no directed cycle: you
        cannot return to your own past.
  \item[A leaf must be incoming.] If $\code{in\_deg}+\code{out\_deg}=1$ (i.e.\ the
        vertex is a degree-1 leaf) we require $\code{in\_deg}=1$: the external leg
        carries the field \emph{in}; the arrow points \emph{into} the leaf.
  \item[No multi-line pure sink.] $\code{out\_deg}=0$ with $\code{in\_deg}\ge 2$ is
        forbidden --- a pure sink absorbing two or more lines is not an admissible
        \msrjd{} vertex.
  \item[No pass-through degree-2 vertex.] $\code{in\_deg}=1$ and
        $\code{out\_deg}=1$ is forbidden --- a one-in/one-out degree-2 vertex is
        again a trivial line-subdivision.
  \item[No two adjacent sources.] A \emph{source} is a vertex with
        $\code{in\_deg}=0$ (and not a leaf). Two adjacent sources are forbidden:
        sources represent noise-injection points, and two of them cannot be wired
        directly together.
\end{description}

This is the \term{source / interaction / leaf} trichotomy that drives all the
downstream physics: \textbf{sources} (in-degree 0, not a leaf) are
\emph{noise-kernel} insertion points; \textbf{interaction vertices} (both in- and
out-edges) are \emph{action-vertex} insertion points; \textbf{leaves} are
\emph{external legs}. The next chapter and \code{degree\_scan.py} re-derive
exactly this trichotomy.

\subsection{Enumerating and deduplicating orientations}

The brute-force orientation sweep is \code{enumerate\_orientations(G, leaves)}
(\file{loop\_diagram\_enumeration.py:295}):

\begin{lstlisting}
def enumerate_orientations(G, leaves):
    edges = list(G.edges(labels=False))
    return [orient_edges(G, [(bits >> i) & 1 for i in range(len(edges))])
            for bits in range(2 ** len(edges))
            if check_orientation_constraints(
                orient_edges(G, [(bits >> i) & 1 for i in range(len(edges))]),
                leaves)]
\end{lstlisting}

It iterates \code{bits} over $0 \ldots 2^{|E|}-1$, decoding each integer into a
length-$|E|$ orientation vector via the bit-extraction \code{(bits >> i) \& 1}
(shift right by $i$, mask the low bit), builds the \code{DiGraph}, and keeps it
iff the constraints pass.

\begin{gotcha}
\label{gotcha:double-build}
\textbf{\code{enumerate\_orientations} builds each candidate \code{DiGraph}
twice} (\file{loop\_diagram\_enumeration.py:297--301}) --- once \emph{inside} the
\code{check\_orientation\_constraints(...)} call, and once again to materialize
the kept result in the list comprehension. This is functionally correct but
doubles the orientation-construction cost; a clean rewrite would build once and
reuse.
\end{gotcha}

The oriented candidates are deduplicated by \code{remove\_isomorphic\_directed}
(\file{loop\_diagram\_enumeration.py:304}) --- the same $O(n^2)$ pattern as the
undirected case, but using the \emph{directed} isomorphism predicate
\code{directed\_graphs\_isomorphic\_with\_labels}, in which arrows must match too.
Each result is carried as the final 4-tuple \code{(D, G, leaves, internal)}.

\section{Graph isomorphism and nauty, from scratch}
\label{sec:iso}

The dedup steps in both stages rest on \emph{graph isomorphism}, which deserves a
proper introduction since it is where almost all the compute time goes at high
order.

\subsection{Why dedup is needed, and what isomorphism means}

After adding edges, and again after orienting, the same \emph{abstract} graph is
produced many times with different vertex \emph{labels} --- ``vertex $5$ here is
vertex $7$ there.'' Two graphs are \term{isomorphic} if one can be turned into
the other purely by \emph{relabeling vertices}: a bijection of vertices that maps
edges to edges. For Feynman-diagram bookkeeping, isomorphic graphs are the
\emph{same diagram} and must be counted exactly once.

Deciding isomorphism naively means trying all $n!$ relabelings --- hopeless even
for modest $n$. The classic engine that does it efficiently is \term{nauty}
(``No AUTomorphisms, Yes?''), by Brendan~McKay. Nauty computes a \emph{canonical
form}: a deterministic relabeling such that two graphs are isomorphic \emph{iff}
their canonical forms are \emph{literally identical}. It does this by iteratively
\emph{refining a coloring} of the vertices --- partitioning them by degree, then
by their neighbors' degrees, and so on --- and backtracking through the residual
symmetries, pruning the search with the automorphisms it discovers along the way.

\begin{note}
You will not see the string \code{nauty} anywhere in this code --- it lives
\emph{inside} SageMath. Every \code{is\_isomorphic} call routes to nauty's
canonical-form machinery. That is the \emph{only} graph-isomorphism technology
this subsystem relies on, and it is what makes deduplication tractable at all.
\end{note}

\subsection{Colored isomorphism: leaves are not interchangeable}

A plain isomorphism may map \emph{any} vertex to any other. But in our diagrams
that is too permissive: \textbf{a leaf must not be mapped onto an internal
vertex}. An external leg and an interaction point play different physical roles,
so a relabeling that swaps the two is \emph{not} a legal symmetry.

The fix is \term{vertex coloring}. Assign color $0$ to leaves and color $1$ to
internal vertices, and ask for a \emph{color-preserving} isomorphism: nauty/Sage
honors the coloring as an initial vertex partition and only searches over
relabelings that keep colors fixed. The undirected predicate
\code{graphs\_isomorphic\_with\_labels} (\file{loop\_diagram\_enumeration.py:84})
implements exactly this:

\begin{lstlisting}
def graphs_isomorphic_with_labels(G1, leaves1, G2, leaves2):
    if G1.order() != G2.order() or G1.size() != G2.size():
        return False
    if len(leaves1) != len(leaves2):
        return False
    leaves1_set = set(leaves1)
    leaves2_set = set(leaves2)
    G1c = G1.copy(); G2c = G2.copy()
    for v in G1.vertices():
        G1c.set_vertex(v, 0 if v in leaves1_set else 1)
    for v in G2.vertices():
        G2c.set_vertex(v, 0 if v in leaves2_set else 1)
    return G1c.is_isomorphic(G2c)
\end{lstlisting}

Three things to read here. First, the \emph{fast rejects}: different vertex count
(\code{order}), different edge count (\code{size}), or different leaf count
returns \code{False} immediately, without invoking nauty. These cheap invariants
are what keep the $O(n^2)$ dedup affordable at small orders. Second, the coloring:
\code{set\_vertex(v, 0\_or\_1)} attaches a payload to each vertex that Sage
interprets as its color. Third, the test itself: \code{G1c.is\_isomorphic(G2c)} is
now a color-preserving comparison, so leaves can only map to leaves.

The directed analog \code{directed\_graphs\_isomorphic\_with\_labels}
(\file{loop\_diagram\_enumeration.py:99}) is structurally identical; the only
difference is that the final \code{is\_isomorphic} on \code{DiGraph}s also
requires the arrows to match.

\section{External tools, from first principles}

This subsystem touches SageMath heavily, three Python standard-library modules,
and --- only transitively, through \code{degree\_scan.py} --- the project's own
serialization and \code{FieldTheory}. It does \emph{not} itself import
\code{sympy}, \code{scipy}, \code{numba}, or \code{networkx}. Each tool is
explained from scratch below; the imports at the top of
\file{loop\_diagram\_enumeration.py:22--25} are:

\begin{lstlisting}
from sage.all import *
from itertools import combinations, combinations_with_replacement
from collections import Counter, defaultdict
from concurrent.futures import ThreadPoolExecutor, as_completed
\end{lstlisting}

\subsection{SageMath}
\label{sec:sage}

\term{SageMath} is a large open-source mathematics system built on top of Python.
It bundles dozens of specialized math libraries (Pari, GAP, Singular, nauty,
NetworkX, NumPy, $\ldots$) behind a uniform Python API, and adds its own types
for exact arithmetic, symbolic expressions (the ``Symbolic Ring'' \code{SR}),
and --- most relevant here --- \emph{graph theory}: the \code{Graph}, \code{DiGraph},
and \code{graphs.*} generator classes.

Because Sage replaces some Python built-ins (e.g.\ integer literals become Sage
\code{Integer}s) and needs its bundled libraries on the path, \emph{this module
can only run inside a SageMath kernel/interpreter}. The module docstring says so
plainly --- \code{"Requires a SageMath kernel."}
(\file{loop\_diagram\_enumeration.py:7}) --- and the tests are run via \code{sage
-python -m pytest}. The whole API is pulled in with the conventional wildcard
import \code{from sage.all import *} (\file{loop\_diagram\_enumeration.py:22}),
which brings \code{Graph}, \code{DiGraph}, \code{graphs}, \code{Integer}, and the
rest into scope.

The specific Sage features this code uses:

\begin{itemize}
  \item \code{graphs.trees(n)} (\file{loop\_diagram\_enumeration.py:149}) --- a
        \emph{generator} yielding every non-isomorphic tree on $n$ vertices,
        exactly once each. This is the brute-force engine of the tree stage; Sage
        generates these efficiently so the user never enumerates trees by hand.
  \item \code{Graph(tree, multiedges=True, loops=False)}
        (\file{loop\_diagram\_enumeration.py:174}) --- an undirected graph that
        \emph{allows} multiple edges between the same pair (\code{multiedges=True},
        needed because a loop-edge can double an existing edge to make a bubble)
        but \emph{forbids} self-loops (\code{loops=False}).
  \item \code{DiGraph(multiedges=True, loops=False)}
        (\file{loop\_diagram\_enumeration.py:265}) --- the directed analog, used to
        hold an oriented topology (a prediagram).
  \item \textbf{Inspection methods:} \code{G.vertices()}, \code{G.edges(labels=...)},
        \code{G.degree(v)}, \code{G.neighbors(v)}, \code{G.is\_connected()},
        \code{G.order()} ($=|V|$), \code{G.size()} ($=|E|$), \code{G.add\_edge(u, v)},
        \code{G.copy()}, \code{G.relabel(dict, inplace=False)}.
  \item \code{G.set\_vertex(v, value)} (\file{loop\_diagram\_enumeration.py:93})
        --- attaches an arbitrary payload (here the color $0$/$1$) to a vertex;
        Sage's isomorphism test treats these payloads as a vertex coloring.
  \item \code{G.is\_isomorphic(H)} (\file{loop\_diagram\_enumeration.py:96}) ---
        the \textbf{nauty-backed} isomorphism test; with vertex colors set first
        it becomes the color-preserving test of \S\ref{sec:iso}.
  \item \textbf{DiGraph-specific:} \code{D.add\_vertices(...)}, \code{D.add\_edge(u,
        v, label)}, \code{D.in\_degree(v)}, \code{D.out\_degree(v)},
        \code{D.is\_directed\_acyclic()}, \code{D.neighbors\_out(v)}.
  \item \textbf{Plotting:} \code{graphics\_array}, \code{G.plot(...)},
        \code{.show()}, \code{.save()} --- used only by the \code{\_plot\_*}
        helpers and the \code{show\_*} API for inline notebook display, not by the
        data pipeline.
\end{itemize}

\subsection{\texttt{itertools}}
\label{sec:itertools}

A Python standard-library module of memory-efficient combinatorial iterators. Two
functions are used:

\begin{itemize}
  \item \code{combinations(iterable, r)} yields every $r$-subset (no repeats,
        order ignored). In \code{generate\_edge\_multisets}
        (\file{loop\_diagram\_enumeration.py:169}),
        \code{combinations(all\_vertices, 2)} enumerates every possible edge ---
        every unordered pair of distinct vertices.
  \item \code{combinations\_with\_replacement(iterable, r)} is like
        \code{combinations} but \emph{allows the same element more than once}. At
        \file{loop\_diagram\_enumeration.py:170},
        \code{combinations\_with\_replacement(possible\_edges, ell)} enumerates
        every multiset of $\ell$ edges to add. ``With replacement'' is what permits
        adding the \emph{same} edge twice --- the double-edge that makes a bubble.
\end{itemize}

\subsection{\texttt{collections}}

\begin{itemize}
  \item \code{defaultdict(list)} --- a dict that auto-creates an empty list on
        first access to a missing key. Used in \code{\_plot\_trees}
        (\file{loop\_diagram\_enumeration.py:349}) to bucket trees by their $j$
        value for grouped display.
  \item \code{Counter} is imported (\file{loop\_diagram\_enumeration.py:24}) but
        \textbf{not actually referenced} anywhere in the file --- a harmless dead
        import (see Gotcha~\ref{gotcha:deadimport}).
\end{itemize}

\subsection{\texttt{concurrent.futures}}
\label{sec:parallel}

\begin{itemize}
  \item \code{ThreadPoolExecutor(max\_workers=n)} runs callables on a pool of OS
        \emph{threads} (not processes). Used in \code{\_enumerate\_topologies\_raw}
        (\file{loop\_diagram\_enumeration.py:248}) and
        \code{\_enumerate\_prediagrams\_raw}
        (\file{loop\_diagram\_enumeration.py:332}) to fan the per-tree /
        per-topology work across threads when \code{n\_threads > 1}.
  \item \code{as\_completed(futures)} yields each submitted future \emph{as it
        finishes} (not in submission order), so results are collected as soon as
        they are ready.
\end{itemize}

\begin{gotcha}
\label{gotcha:threads}
\textbf{Thread-only parallelism is intentional and load-bearing for safety.} The
choice of a thread pool over a process/fork pool is deliberate. Per the project's
macOS fork-safety history, \emph{fork-based multiprocessing crashes the user's
machine} when a kernel forks after BLAS/Cocoa/matplotlib initialization. Threads
avoid \code{fork} entirely. The trade-off is the Python GIL --- but Sage's graph
routines spend most of their time in compiled C (nauty, the graph backend), which
\emph{releases} the GIL, so threads still parallelize the heavy isomorphism work.
The default is \code{n\_threads=1} (serial) everywhere; the pool only activates
when the caller opts in. Do \textbf{not} ``optimize'' this to a
\code{ProcessPoolExecutor} or fork --- it can crash the machine.
\end{gotcha}

\section{The public API}
\label{sec:api}

The module exposes a small, regular surface. The \code{enumerate\_*} functions
return data; the \code{show\_*} variants additionally render the graphs inline in
a notebook and return just the count(s). The colors quoted below are documented in
the corresponding \code{show\_*} docstrings.

\begin{itemize}
  \item \code{enumerate\_all(k, ell, n\_threads=1, max\_vertices\_search=50,
        verbose=True)} (\file{loop\_diagram\_enumeration.py:434}) runs the full
        pipeline \emph{once}, sharing the tree list across stages: it passes
        \code{trees\_with\_j=trees} into \code{\_enumerate\_topologies\_raw} so the
        trees are not regenerated. Returns \code{(trees, topologies, prediagrams,
        counts)} where \code{counts = \{n\_trees, n\_topologies, n\_prediagrams\}}.
  \item \code{enumerate\_trees(k, ell, ...)}
        (\file{loop\_diagram\_enumeration.py:478}) --- a thin wrapper over
        \code{generate\_trees\_with\_constraints}; returns \code{(trees, count)}.
  \item \code{enumerate\_topologies(k, ell, ...)}
        (\file{loop\_diagram\_enumeration.py:498}) --- returns \code{(topologies,
        count)}.
  \item \code{enumerate\_prediagrams(k, ell, ...)}
        (\file{loop\_diagram\_enumeration.py:515}) --- runs topologies then
        orientations; returns \code{(prediagrams, count)}. \textbf{This is the
        canonical entry point the rest of the pipeline calls.}
  \item \code{show\_all} / \code{show\_trees} / \code{show\_topologies} /
        \code{show\_prediagrams} (\file{loop\_diagram\_enumeration.py:462},
        \file{:535}, \file{:552}, \file{:569}) --- display each non-empty stage and
        return the count(s). The prediagram plot colors \emph{leaves} black,
        \emph{sources} red, and other internal vertices light blue, recomputing
        sources on the fly as the vertices with \code{in\_degree == 0}.
\end{itemize}

\begin{note}
The shared-tree optimization in \code{enumerate\_all} matters because tree
generation walks \code{graphs.trees(n)} for every $n$ in the search window, which
is not free. Running the three stages through \code{enumerate\_all} (rather than
calling \code{enumerate\_trees}, \code{enumerate\_topologies}, and
\code{enumerate\_prediagrams} separately) generates the trees exactly once and
threads them through.
\end{note}

\section{The degree-scan feedback bridge (build phase~C)}
\label{sec:degree-scan}

The second file, \file{msrjd/enumeration/degree\_scan.py}, is small but
important: it closes the loop between the \emph{topological} world (this chapter)
and the \emph{algebraic} world (the Taylor-expanded action of Phase~1). The
generator does not look at the action --- but the prediagrams it produces may
\emph{demand} interaction monomials of a higher degree than the action was
expanded to. Phase~C detects that and triggers a re-expansion.

Unlike the generator, this file imports no third-party libraries --- only project
internals (\file{degree\_scan.py:13--14}):

\begin{lstlisting}
from msrjd.core.serialize import load_theory, save_theory, reload_model
from msrjd.core.field_theory import FieldTheory
\end{lstlisting}

\subsection{Reading the degree a theory must support}

\code{max\_vertex\_degree(prediagrams)} (\file{degree\_scan.py:17}) returns the
maximum total degree $\code{in\_degree}(v)+\code{out\_degree}(v)$ over all
\emph{non-leaf} vertices across all prediagrams (or $0$ if none):

\begin{lstlisting}
def max_vertex_degree(prediagrams):
    max_deg = 0
    for (D, G, leaves, internal) in prediagrams:
        leaf_set = set(leaves)
        for v in D.vertices():
            if v in leaf_set:
                continue
            deg = D.in_degree(v) + D.out_degree(v)
            if deg > max_deg:
                max_deg = deg
    return max_deg
\end{lstlisting}

This is the \emph{required Taylor order}: a degree-$d$ internal vertex needs a
$d$-leg interaction monomial, which exists only if the action was expanded to
order $d$. Its sibling \code{scan\_source\_vertices(prediagrams)}
(\file{degree\_scan.py:45}) returns the \emph{set of out-degrees} exhibited by
\emph{source} vertices (non-leaf, \code{in\_degree == 0}) --- these out-degrees
tell the noise-kernel machinery which noise arities are needed.

\subsection{Comparing to the saved order, and re-expanding}

\code{check\_taylor\_order(meta, max\_degree)} (\file{degree\_scan.py:70}) is a
one-liner comparison returning \code{(sufficient, current\_order,
required\_order)} with $\code{sufficient} = (\code{meta['taylor\_order']} \ge
\code{max\_degree})$. The driver is \code{ensure\_taylor\_order(theory\_path,
prediagrams, project\_root=None)} (\file{degree\_scan.py:91}):

\begin{lstlisting}
def ensure_taylor_order(theory_path, prediagrams, project_root=None):
    meta, data = load_theory(theory_path)
    max_deg = max_vertex_degree(prediagrams)
    sufficient, current, required = check_taylor_order(meta, max_deg)
    if sufficient:
        return meta, data
    # Re-expand at the required order
    print(f'Re-expanding theory from order {current} to {required}...')
    model = reload_model(meta, project_root=project_root)
    ft = FieldTheory(model, taylor_order=required)
    ft.expand()
    save_theory(
        theory_path, ft,
        propagator_data=None,  # propagator must be recomputed separately
        stationarity=meta.get('stationarity', True),
        model_file=meta.get('model_file'),
        model_var_name=meta.get('model_var_name'),
    )
    print(f'Theory re-expanded and saved. Note: propagator data must be recomputed.')
    return load_theory(theory_path)
\end{lstlisting}

In the common case (\code{sufficient == True}) the function returns the
already-loaded theory immediately and never touches \code{FieldTheory}. Only on
the re-expansion branch does it reload the model definition
(\code{reload\_model}), rebuild a \code{FieldTheory(model,
taylor\_order=required)}, call \code{.expand()}, re-save, and re-read from disk.
The supporting serialization calls are: \code{load\_theory(path)} returning
\code{(meta, data)} and \code{save\_theory(path, ft, ...)} (read/write a
previously-expanded theory, where \code{meta} carries \code{meta['taylor\_order']}
and, on the re-expansion path, \code{meta['stationarity']},
\code{meta['model\_file']}, \code{meta['model\_var\_name']}); and
\code{reload\_model(meta, project\_root=...)} (re-import the model definition from
\code{meta['model\_file']} / \code{meta['model\_var\_name']}).

\begin{gotcha}
\label{gotcha:stale-prop}
\textbf{\code{ensure\_taylor\_order} leaves the propagator stale after
re-expansion.} It passes \code{propagator\_data=None} to \code{save\_theory} and
prints ``propagator data must be recomputed''
(\file{degree\_scan.py:128--133}). A caller that re-expands and then immediately
uses the propagator without recomputing it will get a wrong or missing
propagator. This is by design --- propagator computation is a separate phase ---
but it is an easy footgun.
\end{gotcha}

\section{Data structures and the downstream contract}

The subsystem passes around a small number of tuple shapes; no dataclasses are
defined here (the \code{TypedDiagram} dataclass lives downstream in
\code{type\_assignment.py}).

\begin{description}
  \item[Tree record] \code{(tree, j, num\_leaves)} --- \code{tree} a Sage
        \code{Graph} (acyclic); \code{j} the leaf surplus; \code{num\_leaves =
        k + j} the leaf count.
  \item[Topology record] \code{(G, leaves, internal)} --- \code{G} a Sage
        \code{Graph} (undirected, multi-edges allowed, exactly $\ell$ loops),
        relabeled leaves-first; \code{leaves} the vertex IDs $0\ldots k-1$;
        \code{internal} the remaining IDs.
  \item[Prediagram record] \code{(D, G, leaves, internal)} --- \emph{the} output
        object. \code{D} a Sage \code{DiGraph} (the oriented graph, edge labels
        carrying the original edge index for multi-edge identity); \code{G} the
        underlying undirected topology; \code{leaves} / \code{internal} as above.
        This 4-tuple is exactly what
        \code{type\_assignment.enumerate\_typed\_diagrams} unpacks and what
        \code{degree\_scan.*} iterate over.
  \item[Counts dict] \code{\{'n\_trees', 'n\_topologies', 'n\_prediagrams'\}},
        built in \code{enumerate\_all} (\file{loop\_diagram\_enumeration.py:454}).
  \item[Edge-multiset] a tuple of $\ell$ vertex-pairs, e.g.\ \code{((0,3),(0,3))}
        for a double-edge bubble between vertices $0$ and $3$.
  \item[Orientation vector] a length-$|E|$ list of bits ($0$ = keep $(u,v)$,
        $1$ = flip to $(v,u)$), produced by \code{(bits >> i) \& 1}.
\end{description}

The end-to-end flow of a single \code{enumerate\_prediagrams(k, ell)} call is then:

\begin{lstlisting}[style=console]
  (k, ell)
     |  generate_trees_with_constraints
  [ (tree, j, num_leaves), ... ]                      trees passing degree bounds
     |  process_tree_parallel  (add ell edges; demand L==ell; constraints;
     |                          relabel leaves-first)
  [ (G, leaves, internal), ... ]  (with isomorphic duplicates)
     |  _remove_isomorphic_undirected   (nauty, leaf-colored)
  [ (G, leaves, internal), ... ]  unique TOPOLOGIES
     |  enumerate_orientations  (2^|E| orientations, keep admissible)
  [ (D, G, leaves, internal), ... ]  (with isomorphic duplicates)
     |  remove_isomorphic_directed   (nauty, leaf-colored, directed)
  [ (D, G, leaves, internal), ... ]  unique PREDIAGRAMS  --->  downstream
\end{lstlisting}

\section{A worked walk-through: \texorpdfstring{$(k,\ell)=(2,1)$}{(k,l)=(2,1)}}

It helps to trace the smallest non-trivial case --- the one-loop correction to a
two-point function. We have $k=2$, $\ell=1$, so by \eqref{eq:jmax} the surplus
runs $j\in\{0,1\}$.

\begin{itemize}
  \item For $j=0$: $\code{num\_leaves}=2$, $\code{v3\_max}=k+j-2=0$ (no branching
        vertices allowed), $\code{v2\_max}=k+3\ell-j-1=4$. The trees here are
        simple paths with two leaves and a short chain of degree-2 vertices
        between them.
  \item For $j=1$: $\code{num\_leaves}=3$, $\code{v3\_max}=1$ (one branching
        vertex allowed --- a ``Y'' joint), $\code{v2\_max}=3$. The third leaf is
        the one the single loop-edge will consume.
\end{itemize}

Each surviving tree then has $\ell=1$ extra edge added in every
\code{combinations\_with\_replacement} way (including the double-edge that makes a
bubble), and only the results with loop number exactly $1$ that pass the
structural constraints are kept. After leaf-colored deduplication this yields
\textbf{9 topologies} for $(k,\ell)=(2,1)$. Orienting those 9 topologies under the
causality constraints of \S\ref{sec:iso} and deduplicating again gives the final
prediagram list that flows downstream to type assignment. (The corresponding
deduplicated counts at the next orders are $(3,1)\to 67$ and $(2,2)\to 289$.)

\section{Performance notes and known limitations}

We collect the sharp edges that were flagged inline, plus the remaining
performance caveats, so they live in one place.

\begin{gotcha}
\label{gotcha:deadimport}
\code{Counter} is imported but unused (\file{loop\_diagram\_enumeration.py:24}) ---
only \code{defaultdict} is referenced. Harmless dead import.
\end{gotcha}

\begin{gotcha}
\label{gotcha:on2}
\textbf{$O(n^2)$ isomorphism dedup.} Both \code{\_remove\_isomorphic\_undirected}
and \code{remove\_isomorphic\_directed} compare each new candidate against
\emph{every} kept representative. The cheap invariant fast-rejects (order, size,
leaf-count) keep this affordable at small orders, but the comparison count is
quadratic in the candidate count and will dominate at high order. A canonical-form
bucket (Sage's \code{Graph.canonical\_label}, hashing the canonical form and
bucketing) would make it near-linear; that is not done here.
\end{gotcha}

\begin{gotcha}
\label{gotcha:edgelabels}
\textbf{Edge labels carry the original edge index.} \code{orient\_edges} adds each
arrow as \code{D.add\_edge(u, v, i)} with \code{i} the edge index
(\file{loop\_diagram\_enumeration.py:269--271}). This is what preserves multi-edge
identity through orientation and into the prediagram \code{D}. Downstream
multi-edge handling (\code{type\_assignment} iterating \code{D.edges()} as 3-tuples
\code{(u, v, label)}) depends on these labels being present and distinct.
\end{gotcha}

\begin{gotcha}
\label{gotcha:leaffirst}
\textbf{The leaf-first labeling is assumed downstream.} \code{relabel\_leaves\_first}
is applied to every accepted topology inside \code{process\_tree\_parallel}
(\file{loop\_diagram\_enumeration.py:213}), and downstream code
(\code{type\_assignment}, \code{degree\_scan}) treats \code{leaves} as the
canonical external-leg set, relying on the relabeling having been done. Any code
path that builds prediagrams \emph{without} going through
\code{process\_tree\_parallel} (e.g.\ hand-built mocks in tests) will not have
leaf-first labels.
\end{gotcha}

\begin{gotcha}
\label{gotcha:sentinel}
\code{count\_cycles\_sage} returns $-1$ for a disconnected graph
(\file{loop\_diagram\_enumeration.py:69--71}), which is not $\ell$ for any
$\ell\ge 0$; the disconnected graph is therefore rejected automatically by the
\code{== ell} test in \code{process\_tree\_parallel}. The $-1$ is a deliberate
sentinel, not an error.
\end{gotcha}

The remaining structural limitations were stated in context and are recalled here
for completeness: the \code{+10} vertex-search cap (Gotcha~\ref{gotcha:plus10})
can in principle undercut the proven completeness bound at high order; the
double-build in \code{enumerate\_orientations} (Gotcha~\ref{gotcha:double-build})
doubles the orientation-construction cost; the $\ell=0$ constraint skip
(Gotcha~\ref{gotcha:ell0-skip}) means the topology constraint set differs at tree
level; and the whole module requires a SageMath kernel (\S\ref{sec:sage}). None of
these affect correctness at the small orders the pipeline routinely runs, but all
are worth keeping in view when pushing $\ell$ to three loops and beyond.

\section{Where to go next}

The prediagram list this chapter produces is consumed in
Chapter~\ref{ch:type-assignment}, where each bare \code{(D, G, leaves, internal)}
tuple is walked and every consistent assignment of propagators to edges and
interaction monomials to vertices is tried, the invalid ones discarded, and a
fully-typed \code{TypedDiagram} emitted for each survivor. That is the step where
the physics --- the response-versus-correlation propagator distinction, the
specific vertices the action provides, and the symmetry factor $\Sym(\Gamma)$ ---
finally re-enters. The completeness guaranteed here (every shape) and the validity
guaranteed there (only labelings the action supports) together ensure the diagram
sum the engine evaluates is exactly the right one.
